\section{Roles and Protocol}
\label{sec:protocol}

\subsection{Roles}

TRACE instantiates the three RATS roles~\citep{rats2021}. The \emph{Attester} is the
agent runtime: it holds the TEE-bound signing key, gathers Evidence (the platform
measurement chain), and issues a Trust Record. The \emph{Verifier} appraises that
Evidence against reference values and policy and emits an appraisal (the record's
\texttt{appraisal} claim, in the EAR format). The \emph{Relying Party} (a
downstream service, an auditor, or a partner) consumes the record and the
appraisal to make a trust decision.

TRACE follows the RATS \emph{passport} model: the Attester obtains its evidence and
ships a self-contained Trust Record that travels with the workload. Any party
verifies it offline, with no callback to the issuer and no vendor in the trust path
beyond the silicon root and the transparency-log operator
(Section~\ref{sec:verification}).

\begin{figure}[h]
\centering
\begin{tikzpicture}[
  font=\footnotesize,
  >={Stealth[round]},
  role/.style={draw, rounded corners, align=center, minimum width=2.9cm, minimum height=2.6em},
  note/.style={font=\scriptsize\itshape, text=black!70},
]
\node[role, fill=blue!5] (att) {Attester\\(agent runtime / TEE)};
\node[role, fill=gray!8, right=2.6cm of att] (ver) {Verifier};
\node[role, fill=green!6, below=1.1cm of att] (rp) {Relying Party};

\draw[->] (att) -- node[note, above]{Evidence} (ver);
\draw[->] (ver) to[out=250,in=20] node[note, right, pos=0.4]{appraisal (EAR)} (rp);
\draw[->] (att) -- node[note, left]{Trust Record} (rp);
\node[note, below=0.2cm of rp] {verifies offline, no issuer callback};
\end{tikzpicture}
\caption{TRACE in the RATS passport model. The Attester ships a self-contained
Trust Record; the Verifier's appraisal travels in the record so the Relying Party
appraises it without contacting the issuer.}
\label{fig:roles}
\end{figure}

\subsection{Key hierarchy}

Signing keys form a chain rooted in silicon:
silicon root key (vendor-managed, hardware-bound) $\rightarrow$ platform attestation
key (Intel TDX Quote-signing key, AMD VCEK/VLEK, NVIDIA NRAS) $\rightarrow$ workload
attestation key (TEE-bound, ephemeral) $\rightarrow$ record-signing key (per
workload, optionally per session). The record-signing key is the one named in
\texttt{cnf}; workload-level keys SHOULD rotate at TEE-image boundaries, and
verifiers MUST consult current revocation status at verification time.

\subsection{Single-organization flow (Phase 1)}

\begin{enumerate}[leftmargin=1.5em]
  \item \textbf{Startup attestation.} The runtime launches inside a TEE, derives its
    record-signing key, and obtains a platform attestation report binding that key.
  \item \textbf{Issuance.} For each governed execution (an invocation, or a session),
    the runtime assembles a Trust Record (\texttt{model}, \texttt{runtime},
    \texttt{policy}, \texttt{data\_class}, \texttt{tool\_transcript},
    \texttt{build\_provenance}), signs it (Section~\ref{sec:model:sig}), and
    optionally registers it on the transparency log, filling \texttt{transparency}.
  \item \textbf{Appraisal and consumption.} A Verifier appraises the record and a
    Relying Party acts on it, both offline.
\end{enumerate}

Issuance granularity is profile-defined: the cMCP profile issues one record per
session, at session close, binding that session's audit chain; a profile that needs
finer evidence may issue per invocation. In all cases the \texttt{subject} is the
workload's SPIFFE identity~\citep{spiffe2018}, bound to the TEE measurement so the
record names which agent acted.

\subsection{Cross-organization flow (Phase 2)}

When two runtimes in different trust domains participate (for example a gateway
and an upstream server, each in its own TEE), each is an independent Attester. Each
produces its own Trust Record, signed by its own \texttt{cnf} key and bound to its
own platform attestation, and both carry the \emph{same} \texttt{subject}/session
identifier. The shared identifier links the two records without sharing private key
material; a Verifier confirms the linkage and appraises each record independently, so
tampering with one does not affect the other. cMCP~\citep{cmcp2026} demonstrates this
as a dual-TEE protocol at the MCP boundary.

\begin{figure}[h]
\centering
\begin{tikzpicture}[
  font=\footnotesize,
  >={Stealth[round]},
  tee/.style={draw, dashed, thick, rounded corners, align=center, minimum width=3.6cm, minimum height=3.0em, inner sep=6pt},
  rec/.style={draw, rounded corners, fill=blue!5, align=center, minimum height=1.9em},
  note/.style={font=\scriptsize\itshape, text=black!70},
]
\node[tee, draw=blue!55!black] (a) {\textbf{Org A: Gateway TEE}\\ \texttt{cnf}: $K_A$};
\node[tee, draw=orange!75!black, right=3.0cm of a] (b) {\textbf{Org B: Server TEE}\\ \texttt{cnf}: $K_B$};
\draw[<->, thick] (a.north) to[out=60,in=120] node[note, above]{shared session id $\sigma$} (b.north);
\node[rec, below=0.8cm of a] (ra) {Trust Record\\(signed by $K_A$)};
\node[rec, below=0.8cm of b] (rb) {Trust Record\\(signed by $K_B$)};
\draw[->] (a) -- (ra);
\draw[->] (b) -- (rb);
\node[draw, rounded corners, fill=green!6, align=center, minimum height=1.9em,
      below=0.8cm of $(ra)!0.5!(rb)$] (v) {Verifier: links via $\sigma$, appraises each independently};
\draw[->] (ra.south) -- (v.north west);
\draw[->] (rb.south) -- (v.north east);
\end{tikzpicture}
\caption{Cross-organization attestation. Two independent Attesters issue separately
key-bound Trust Records carrying a shared session identifier; the Verifier links them
without either side sharing key material.}
\label{fig:crossorg}
\end{figure}
